% ===================== CHAPITRE 3 — CONCEPTION =======================
\chapter{Conception}

\section*{Introduction}
Ce chapitre traduit l'analyse en une \textbf{conception} détaillée du système. Il
présente d'abord l'architecture (globale, physique et logique), puis le modèle
statique du domaine (diagramme de classes) et enfin les scénarios dynamiques des
fonctionnalités principales (diagrammes de séquence, d'états-transitions et
d'activité).

% ----------------------------------------------------------------------
\section{Architecture du système}
L'architecture définit l'organisation des composants du système et la manière dont
ils communiquent. Elle est présentée à trois niveaux : globale, physique et
logique.

\subsection{Architecture globale}
L'architecture globale (figure~\ref{fig:archi_globale}) offre une vue d'ensemble
des composants. Un poste client (navigateur) dialogue en \texttt{HTTP/HTTPS} avec
l'application web déployée sur \textbf{WildFly}, structurée en couches
(présentation, contrôle, métier, persistance) ; celle-ci accède à la base
\textbf{Oracle} en JDBC. La base alimente, d'une part, \textbf{Power BI} via les
vues \texttt{V\_BI\_*} et, d'autre part, la \textbf{chaîne décisionnelle Python}
(ETL puis segmentation).

\begin{figure}[htbp]
  \centering
  \includegraphics[width=0.92\textwidth]{archi_globale.png}
  \caption{Architecture globale du système}
  \label{fig:archi_globale}
\end{figure}

\subsection{Architecture physique}
L'architecture physique (figure~\ref{fig:archi_physique}) représente le
\textbf{déploiement} des composants sur les nœuds matériels selon une organisation
\textbf{3-tiers} : le poste client, le serveur d'applications WildFly hébergeant
\texttt{gestion-agences.war} (accès à la base via la source de données JTA
\texttt{OracleDS}), le serveur de base de données Oracle (\texttt{FREEPDB1}) et le
poste décisionnel exécutant Power BI.

\begin{figure}[htbp]
  \centering
  \includegraphics[width=0.85\textwidth]{archi_physique.png}
  \caption{Architecture physique (diagramme de déploiement 3-tiers)}
  \label{fig:archi_physique}
\end{figure}

\subsection{Architecture logique}
L'architecture logique (figure~\ref{fig:archi_logique}) décompose le système en
\textbf{couches} faiblement couplées suivant le patron \textbf{MVC} : la couche
\emph{présentation} (JSP/JSTL, ApexCharts), la couche \emph{contrôle} (Servlets et
le filtre \texttt{NotificationFilter}), la couche \emph{métier} (services EJB) et
la couche \emph{persistance} (entités JPA et Hibernate dialoguant avec Oracle via
l'\texttt{EntityManager}).

\begin{figure}[htbp]
  \centering
  \includegraphics[width=0.62\textwidth]{archi_logique.png}
  \caption{Architecture logique en couches de l'application}
  \label{fig:archi_logique}
\end{figure}

% ----------------------------------------------------------------------
\section{Diagramme de classes}
La figure~\ref{fig:class} présente le \textbf{diagramme de classes} du modèle
métier, qui représente les entités persistantes du domaine, leurs attributs et
les associations qui les relient.

\begin{figure}[p]
  \centering
  \rotatebox{90}{%
    \includegraphics[width=0.92\textheight,height=\textwidth,keepaspectratio]{class.png}}
  \caption{Diagramme de classes du modèle métier}
  \label{fig:class}
\end{figure}

L'entité \textbf{\texttt{Utilisateur}} occupe une position centrale. Le
rattachement d'un \texttt{Utilisateur}, d'un \texttt{Client} ou d'un
\texttt{Objectif} à son \texttt{Agence}, ainsi que d'un \texttt{Client} à son
\texttt{Gestionnaire}, est matérialisé dans le code par des \textbf{clés
étrangères} (\texttt{SK\_AGENCE}, \texttt{SK\_GESTIONNAIRE}). Les entités de
demande (\texttt{DemandeClient}, \texttt{DemandeEmploye},
\texttt{DemandeModification}, \texttt{DemandeAgence}), ainsi que
\texttt{JournalAction}, \texttt{Notification} et \texttt{Pointage}, référencent un
\texttt{Utilisateur} par une association JPA \texttt{@ManyToOne} ;
\texttt{DemandeAgence} possède en outre une association \texttt{@ManyToMany} avec
\texttt{Utilisateur} (\texttt{utilisateursDemandes}). Enfin, \texttt{SecurityUser}
est relié à \texttt{Utilisateur} par une dépendance d'authentification. Le modèle
ne comporte ni héritage ni composition : les liens sont des associations, dont
certaines sont implémentées par des clés étrangères et d'autres par des relations
JPA explicites.

% ----------------------------------------------------------------------
\section{Modèle physique de données}
Au modèle de classes correspond, au niveau de la base de données, un
\textbf{schéma relationnel} : chaque entité persistante est traduite en une table
Oracle, chaque attribut en une colonne et chaque association en une \textbf{clé
étrangère}. La figure~\ref{fig:er} présente ce schéma relationnel, obtenu par
\emph{mapping} objet-relationnel (JPA/Hibernate) à partir des entités du domaine.

\begin{figure}[htbp]
  \centering
  \includegraphics[width=\textwidth]{er.png}
  \caption{Schéma relationnel de la base de données}
  \label{fig:er}
\end{figure}

La table \texttt{AGENCE} constitue le référentiel central autour duquel gravitent
les employés (\texttt{B\_UTILISATEURS}), les clients (\texttt{B\_CLIENTS}) et les
objectifs (\texttt{B\_OBJECTIF}), reliés par la clé \texttt{SK\_AGENCE}. La table
\texttt{B\_UTILISATEURS} est à son tour référencée par les pointages, les demandes
et le journal d'audit. Le tableau~\ref{tab:tables} décrit le rôle des principales
tables du schéma.

\begin{table}[htbp]
\centering
\caption{Description des principales tables de la base}
\label{tab:tables}
\renewcommand{\arraystretch}{1.2}
\begin{tabularx}{\textwidth}{|>{\bfseries\ttfamily}p{3.6cm}|X|p{2.8cm}|}
\hline
\textmd{\textbf{Table}} & \textbf{Rôle} & \textbf{Clé primaire}\\
\hline
AGENCE & Référentiel des agences du réseau (libellé, adresse, district). & \texttt{SK\_AGENCE}\\
\hline
B\_UTILISATEURS & Employés et utilisateurs de l'application. & \texttt{SK\_UTILISATEUR}\\
\hline
B\_CLIENTS & Portefeuille clients (type, statut, secteur, adresse). & \texttt{SK\_CLIENT}\\
\hline
B\_OBJECTIF & Objectifs commerciaux et production par période. & \texttt{ID\_OBJECTIF}\\
\hline
POINTAGE & Pointages de présence des employés. & \texttt{ID\_POINTAGE}\\
\hline
GESTIONNAIRE & Gestionnaires et rattachement de portefeuille. & \texttt{SK\_UTILISATEUR}\\
\hline
SECURITY\_USERS & Comptes, mots de passe et rôles (sécurité). & \texttt{ID}\\
\hline
DEMANDE\_CLIENT & Demandes de création de clients. & \texttt{ID\_DEMANDE}\\
\hline
DEMANDE\_EMPLOYE & Demandes de création d'employés. & \texttt{ID\_DEMANDE}\\
\hline
DEMANDES\_MODIF. & Demandes de modification de données. & \texttt{ID}\\
\hline
JOURNAL\_ACTION & Journal d'audit des actions. & \texttt{ID\_JOURNAL}\\
\hline
NOTIFICATION & Notifications adressées aux utilisateurs. & \texttt{ID}\\
\hline
\end{tabularx}
\end{table}

% ----------------------------------------------------------------------
\section{Sécurité et contrôle d'accès}
La sécurité de l'application repose sur trois mécanismes complémentaires :
l'authentification, la gestion des rôles et le filtrage des requêtes.

\subsection{Authentification et session}
Les comptes sont stockés dans la table \texttt{SECURITY\_USERS} (identifiant, mot
de passe et rôle). Le \texttt{LoginServlet} délègue la vérification au
\texttt{SecurityService} ; en cas de succès, une \texttt{HttpSession} est ouverte
et l'\texttt{Utilisateur} métier correspondant y est associé. Le filtre
\texttt{NotificationFilter} intercepte ensuite les requêtes pour vérifier la
présence d'une session valide avant tout accès aux pages protégées, et redirige
vers la page de connexion dans le cas contraire.

\subsection{Rôles et matrice des droits}
Trois rôles structurent les droits d'accès. Le tableau~\ref{tab:permissions}
présente la matrice des permissions par fonctionnalité.

\begin{table}[htbp]
\centering
\caption{Matrice des droits d'accès par rôle}
\label{tab:permissions}
\renewcommand{\arraystretch}{1.25}
\begin{tabularx}{\textwidth}{|X|c|c|c|}
\hline
\textbf{Fonctionnalité} & \textbf{ADMIN} & \textbf{DIR. COMM.} & \textbf{USER}\\
\hline
Gestion des agences, employés, clients & Oui & Non & Non\\
\hline
Validation des demandes & Oui & Modif. employés & Non\\
\hline
Soumission de demandes & Oui & Oui & Ses informations\\
\hline
Consultation des objectifs et indicateurs & Oui & Oui & Autorisé\\
\hline
Pointage de présence & Oui & Oui & Oui\\
\hline
Supervision du volet décisionnel & Oui & Non & Non\\
\hline
Consultation du journal d'audit & Oui & Non & Non\\
\hline
\end{tabularx}
\end{table}

Ce cloisonnement garantit qu'un utilisateur ne peut accéder qu'aux fonctionnalités
autorisées par son rôle : une tentative d'accès direct à une ressource réservée
est bloquée par le contrôle de session et de rôle.

% ----------------------------------------------------------------------
\section{Diagrammes de séquence}
Les diagrammes de séquence décrivent le déroulement dynamique des principales
fonctionnalités, en montrant les échanges entre la couche de présentation, les
servlets, les services métier et la base de données.

\subsection{S'authentifier}
La figure~\ref{fig:seq_auth} décrit le déroulement dynamique de l'authentification.

\begin{figure}[htbp]
  \centering
  \includegraphics[width=0.92\textwidth]{seq_auth.png}
  \caption{Diagramme de séquence -- Authentification}
  \label{fig:seq_auth}
\end{figure}

Le \texttt{LoginServlet} reçoit la requête \texttt{POST /login} et délègue la
vérification au \texttt{SecurityService}, qui interroge la table
\texttt{SECURITY\_USERS} via l'\texttt{EntityManager}. Le fragment combiné
\texttt{alt} distingue les deux issues : en cas de succès, l'\texttt{Utilisateur}
métier est associé (ou créé), une session est ouverte et l'utilisateur est
redirigé vers le tableau de bord ; en cas d'échec, un message d'erreur est
renvoyé.

\subsection{Créer une agence}
La figure~\ref{fig:seq_agence} décrit le scénario de création d'une agence.

\begin{figure}[htbp]
  \centering
  \includegraphics[width=0.92\textwidth]{seq_agence.png}
  \caption{Diagramme de séquence -- Création d'une agence}
  \label{fig:seq_agence}
\end{figure}

L'\texttt{AgenceServlet} reçoit la requête de création et délègue à
l'\texttt{AgenceService}, qui persiste la nouvelle \texttt{Agence} via
l'\texttt{EntityManager} (\texttt{INSERT}), puis recharge la liste des agences
renvoyée à la page \texttt{agences.jsp}. Le même schéma s'applique aux opérations
de modification et de suppression.

\subsection{Gérer un client}
La figure~\ref{fig:seq_client} décrit l'enregistrement et l'affectation d'un
client.

\begin{figure}[htbp]
  \centering
  \includegraphics[width=0.92\textwidth]{seq_client.png}
  \caption{Diagramme de séquence -- Gestion d'un client}
  \label{fig:seq_client}
\end{figure}

Le \texttt{ClientServlet} transmet la demande au \texttt{ClientService}, qui
recherche le \texttt{Gestionnaire} concerné, associe le \texttt{Client} à ce
dernier, puis le persiste via l'\texttt{EntityManager}. Le portefeuille mis à jour
est renvoyé à la page \texttt{client.jsp}.

\subsection{Pointer un employé}
La figure~\ref{fig:seq_pointage} décrit le scénario de pointage.

\begin{figure}[htbp]
  \centering
  \includegraphics[width=0.95\textwidth]{seq_pointage.png}
  \caption{Diagramme de séquence -- Pointage d'un employé}
  \label{fig:seq_pointage}
\end{figure}

Le \texttt{PointageServlet} transmet l'action au \texttt{PointageService}, qui
recherche le pointage du jour. Le fragment \texttt{alt} distingue l'\texttt{INSERT}
(premier pointage, statut \texttt{PRESENT}/\texttt{RETARD} au-delà de 9 h) de
l'\texttt{UPDATE} (pointage déjà existant).

\subsection{Valider une demande}
La figure~\ref{fig:seq_demande} détaille le traitement d'une demande par
l'administrateur.

\begin{figure}[htbp]
  \centering
  \includegraphics[width=0.92\textwidth]{seq_demande.png}
  \caption{Diagramme de séquence -- Validation d'une demande}
  \label{fig:seq_demande}
\end{figure}

Le \texttt{GestionModificationsServlet} transmet la décision au service métier,
qui charge la demande, met à jour son statut, persiste la modification puis
journalise l'action via le \texttt{JournalService}.

% ----------------------------------------------------------------------
\section{Cycle de vie d'une demande}
Les circuits de validation constituent un mécanisme transverse important. Leur
comportement est modélisé par un diagramme d'états-transitions et un diagramme
d'activité.

\subsection{Diagramme d'états-transitions}
La figure~\ref{fig:state_demande} modélise le cycle de vie d'une demande.

\begin{figure}[htbp]
  \centering
  \includegraphics[width=0.6\textwidth]{state_demande.png}
  \caption{Diagramme d'états-transitions -- Cycle de vie d'une demande}
  \label{fig:state_demande}
\end{figure}

Une demande naît à l'état \texttt{EN\_ATTENTE}, passe en \texttt{EN\_COURS} lors de
sa prise en charge, puis atteint un état final \texttt{EFFECTUE} ou
\texttt{REFUSE}, en cohérence avec la contrainte de la table des demandes.

\subsection{Diagramme d'activité}
La figure~\ref{fig:activity_demande} représente le processus métier de traitement
d'une demande.

\begin{figure}[htbp]
  \centering
  \includegraphics[width=0.6\textwidth]{activity_demande.png}
  \caption{Diagramme d'activité -- Traitement d'une demande}
  \label{fig:activity_demande}
\end{figure}

Le directeur soumet une demande, notifiée à l'administrateur. Le nœud de décision
oriente le flux : en cas de validation, le statut passe à \texttt{EFFECTUE} et la
modification est appliquée ; en cas de refus, le statut passe à \texttt{REFUSE}
avec un commentaire. Les deux branches convergent vers la journalisation et la
notification du demandeur.

\section*{Conclusion}
Ce chapitre a défini l'architecture du système, son modèle de classes, son schéma
relationnel, ses mécanismes de sécurité et les scénarios dynamiques de ses
fonctionnalités clés. Le chapitre suivant présente le \textbf{volet décisionnel} :
la chaîne ETL, le modèle en étoile, les tableaux de bord et la segmentation des
agences.
